\section{Casos de estudio}
\label{sec:case-studies}

Las cifras agregadas (\S\ref{sec:evaluation}) muestran \emph{que} la
tuber\'ia rinde bien; los cuatro casos siguientes muestran \emph{por
qu\'e}: instancias concretas de la disciplina de validaci\'on
encontrando y resolviendo un problema real, cada una ilustrando una
clase de fallo distinta que una evaluaci\'on puramente num\'erica no
sacar\'ia a la luz por s\'i sola.

\subsection{La cadena multietapa de RoningLoader}
\label{sec:case-roning}

La cadena de infecci\'on de RoningLoader, establecida mediante
ingenier\'ia inversa manual, no es un binario sino cuatro: un cargador
que descifra y suelta un driver rootkit en modo kernel usado para
ocultar su propio proceso y ficheros, una DLL independiente cuya
\'unica funci\'on es enumerar y terminar productos de seguridad, y un
cliente de acceso remoto que establece el canal C2 real. La capacidad
de cada componente es real y significativa por s\'i misma; ninguno de
los tres componentes soltados aparece como importaci\'on o cadena en el
propio binario cargador, ya que se descifran en memoria o se sueltan a
disco \'unicamente en tiempo de ejecuci\'on.

La Tabla~\ref{tab:results} ya muestra la consecuencia num\'erica:
puntuar el cargador de forma aislada alcanza una exhaustividad de 0.68
frente a una verdad de referencia escrita para la cadena completa,
porque aproximadamente un tercio de las t\'ecnicas que un analista
humano atribuye correctamente a la familia ``RoningLoader'' solo son
observables analizando directamente los componentes soltados. El
bucle de reenv\'io (\S\ref{sec:resubmission}) existe precisamente
porque esto es una propiedad estructural del malware multietapa, no
una peculiaridad de una muestra, y cierra 16 de esos 32 puntos de
exhaustividad al dar a cada componente su propio an\'alisis
independiente en lugar de intentar inferir su capacidad a partir de la
evidencia del padre.

\subsection{El error de agregaci\'on del instalador}
\label{sec:case-installer}

Al principio de la validaci\'on, el informe est\'atico de una muestra
empaquetada como instalador mostraba evidencia consistente solo con
\emph{uno} de sus varios componentes incluidos, pese a que la
l\'ogica de extracci\'on en varias pasadas (\S\ref{sec:static-rules})
hab\'ia procesado todos ellos. Al trazar directamente el c\'odigo de
generaci\'on del informe, sin asumir que el fallo estaba en la propia
l\'ogica de extracci\'on, se encontr\'o el defecto real una etapa
despu\'es: \code{\_analyze\_installer}, la funci\'on que agrega los
resultados por componente en el informe \'unico de nivel superior de
la muestra, estaba \emph{sustituyendo} silenciosamente el an\'alisis de
un componente hijo por el de otro bajo una condici\'on espec\'ifica, en
lugar de fusionar ambos.

El error se confirm\'o, no solo se infiri\'o del s\'intoma, a trav\'es
del propio campo \code{hash\_match} del informe, un valor registrado
precisamente para este tipo de comprobaci\'on cruzada, que mostraba que
el hash reportado pertenec\'ia a un fichero extra\'ido distinto del que
realmente aparec\'ia en el cuerpo del informe. Corregir la l\'ogica de
agregaci\'on para fusionar los hallazgos de cada componente extra\'ido,
en lugar de solo el procesado m\'as recientemente, restaur\'o la
evidencia faltante y mejor\'o directamente la exhaustividad en la
muestra afectada tras volver a evaluarla.

La lecci\'on se generaliz\'o m\'as all\'a de esta correcci\'on
concreta: motiv\'o tratar cada nueva capacidad de la tuber\'ia, el
bucle de reenv\'io en primer lugar, como necesitada de su propio paso
expl\'icito de corroboraci\'on (manifiestos de linaje, comprobaciones
cruzadas de hash) en lugar de confiar en que ``el c\'odigo se
ejecut\'o sin error'' implique ``el c\'odigo produjo el resultado
correcto''. Encontrar y corregir esta clase de error antes de que
pudiera subestimar en silencio la capacidad real de una muestra es en
s\'i mismo una demostraci\'on concreta de para qu\'e sirve la
disciplina de validaci\'on de la tuber\'ia.

\subsection{El hueco de verdad de referencia de WhiteSnakeStealer}
\label{sec:case-wsnake}

La \S\ref{sec:fp-audit} resume el resultado; el proceso merece
detallarse una vez como ejemplo concreto de lo que significa en la
pr\'actica ``auditado frente al texto completo del informe''. El
resultado de la tuber\'ia para WhiteSnakeStealer inclu\'ia
\technique{T1620} (Reflective Code Loading) con evidencia est\'atica,
llamadas espec\'ificas a la BCL de .NET consistentes con la carga de
ensamblados en memoria. La propia tabla de T\'ecnicas ATT\&CK de la
muestra, sin embargo, no ten\'ia fila de \technique{T1620}, por lo que
la evaluaci\'on autom\'atica inicialmente puntu\'o esto como un falso
positivo.

Al leer la secci\'on narrativa del informe, en lugar de confiar solo en
su tabla resumen, se encontr\'o un p\'arrafo expl\'icito que
describ\'ia exactamente el mecanismo de carga reflexiva que hab\'ia
detectado la regla est\'atica: exist\'ia confirmaci\'on independiente
del analista, simplemente no se hab\'ia trasladado a la tabla que lee
el extractor autom\'atico de verdad de referencia. Esto es un hueco de
documentaci\'on en el informe fuente, no un error de la tuber\'ia, y se
corrigi\'o a\~nadiendo la fila que faltaba al propio informe (y
regenerando la verdad de referencia a partir de \'el), en lugar de
ajustar la confianza de la tuber\'ia o suprimir el hallazgo. Se
presenta aqu\'i espec\'ificamente porque dos casos estructuralmente
similares en RoningLoader (\S\ref{sec:fp-audit}) se resolvieron en el
sentido opuesto, dejados abiertos, precisamente porque les faltaba
este tipo de confirmaci\'on textual independiente. Juntos, los dos
desenlaces muestran que la evidencia propia de la tuber\'ia se trata
igual independientemente de si la discrepancia termin\'o
favoreci\'endola.

\subsection{La cascada de falsos positivos de la clave XOR}
\label{sec:case-xor}

La fuerza bruta XOR de un solo byte (\code{config\_parser.py}) prueba
exhaustivamente las 256 claves posibles contra regiones candidatas de
texto cifrado y compara el resultado con una expresi\'on regular de IP
en formato punteado, la t\'ecnica que recuper\'o la IP de C2 gen\'uina
de RoningLoader (\code{202.95.11.173}) sin conocimiento previo de la
clave (\S\ref{sec:limitations}, punto~7). Ejecutada sin salvaguardas
contra una segunda muestra, WhiteSnakeStealer, la misma rutina
produjo 8 coincidencias de cuarteto punteado sint\'acticamente
v\'alidas pero completamente espurias a partir de una \'unica clave
incorrecta (\code{0x2e}, ASCII \code{.}): una tabla de enteros
repetitiva con relleno de ceros en los datos compilados del binario
result\'o decodificar a ocho direcciones IP falsas bajo esa \'unica
clave, ya que los bytes nulos se decodifican como caracteres \code{.}
literales bajo ella y una estructura repetitiva en los datos de origen
produce una coincidencia falsa repetitiva.

Dos filtros cerraron el hueco, cada uno justificado por una propiedad
distinta que un valor incrustado gen\'uino tiene y un accidente de
decodificaci\'on no. Primero, una \textbf{comprobaci\'on de l\'imite
limpio}: el byte inmediatamente anterior y posterior a una coincidencia
debe ser \code{0x00} o el propio byte de clave, algo que una cadena
real incrustada y terminada en nulo cumple y que una secuencia de
bytes arbitraria decodificada por una clave no relacionada
generalmente no. Segundo, un \textbf{l\'imite de densidad de
coincidencias} de dos coincidencias con forma de IP por clave en
cualquier lugar del fichero, bajo el principio de que un valor de
configuraci\'on gen\'uino aparece una vez, o raramente un pu\~nado de
veces, mientras que una cascada producida al decodificar datos
estructurados no relacionados bajo la clave equivocada no se
autolimita. Verificado despu\'es en ambas direcciones: la IP de C2
gen\'uina de RoningLoader sigue produciendo exactamente una coincidencia
y sobrevive al l\'imite de densidad; las ocho coincidencias espurias de
WhiteSnakeStealer ahora se descartan correctamente.

La lecci\'on se generaliz\'o m\'as all\'a de esta correcci\'on
concreta: es lo que motiv\'o la pol\'itica expl\'icita del proyecto
(\S\ref{sec:limitations}, punto~9) de que cualquier nueva capacidad de
recuperaci\'on basada en b\'usqueda exhaustiva debe validarse frente a
una segunda muestra, no solo la que motiv\'o su construcci\'on, antes
de confiar en ella. De haberse declarado esta capacidad lista para
producci\'on solo con RoningLoader, la tuber\'ia habr\'ia salido con un
generador de falsos positivos que simplemente todav\'ia no se hab\'ia
activado.
